\chapteropen{08}

\begin{cele}
\item rozstrzygnąć, które książki warto notować, a które wystarczy przeczytać, i policzyć koszt czasowy notowania;
\item napisać jednozdaniowe streszczenie książki i zamienić je w trzy konkretne działania;
\item zaplanować powtórki w odstępach oraz przypominanie treści przez rozmowę zamiast ponownego czytania;
\item ocenić zakreślanie, marginalia i cytaty jako narzędzia, a nie jako dowód zrozumienia;
\item prowadzić dziennik czytania i łączyć książki ze sobą, nie trzymając ich w osobnych szufladach.
\end{cele}

\section{Po co w ogóle zapamiętywać}

Większość tego, co czytasz, zanika szybko. Dokładny przebieg tej utraty zależy od tekstu, od tego, jak bardzo był dla ciebie nowy, i od tego, co zrobiłeś z nim później, ale kierunek jest znany od dawna: bez powrotu do treści zostaje niewiele \verify{po kilku tygodniach zwykle niewiele z przeczytanej książki pamięta się tak, żeby dało się z niej korzystać}. Dla kogoś, kto czyta dwie książki w miesiącu, oznacza to, że po roku większość lektur istnieje wyłącznie jako wrażenie.

Ta obserwacja jest powodem, dla którego ludzie zaczynają notować, i zarazem powodem, dla którego notowanie tak często zawodzi. Notuje się to, co łatwo zapisać: cytaty, nagłówki, numery stron. Tymczasem to, co warto zachować, jest zwykle czymś innym: wnioskiem, który zmienia jedną decyzję, albo przykładem, do którego wrócisz za pół roku.

Warto rozdzielić dwa cele, bo wymagają innych narzędzi. Pierwszy to \emph{przechowanie} treści: możesz kiedyś sprawdzić, co autor napisał o danym temacie. Drugi to \emph{użycie}: potrafisz tym coś zrobić w pracy, w rozmowie, w decyzji. Notatki dobrze służą pierwszemu celowi, a drugiemu tylko wtedy, gdy są pisane z myślą o konkretnym zastosowaniu. Większość tego rozdziału dotyczy właśnie tego rozróżnienia.

\section{Kiedy notować, a kiedy tylko czytać}

Nie każda książka wymaga notatek i nie każda notatka jest tyle warta, ile kosztuje. Zanim zaczniesz pisać, zadaj sobie trzy pytania. Czy będę z tej książki \emph{korzystać} w ciągu najbliższych miesięcy? Czy będę do niej \emph{wracać}, na przykład po konkretną technikę albo definicję? Czy ma dla mnie wartość sama lektura, na przykład powieść albo książka, której sens leży w doświadczeniu, a nie w tezie? Jeśli na wszystkie trzy odpowiedź brzmi ,,nie'', wystarczy czytać. Jeśli choć na jedno odpowiadasz ,,tak'', warto poświęcić notatkom trochę czasu.

Koszt notowania trzeba policzyć, bo łatwo go nie zauważyć. Typowa książka ma około \val{typical.words} słów. Przy szybkości 250 słów na minutę, która jest tu przyjętym punktem odniesienia, czytanie zajmuje około \val{typical.hours} godziny. Dwanaście takich książek w roku to około \val{need.12} minuty dziennie, a pięćdziesiąt dwie to około \val{need.52} minuty. Ten rachunek nie uwzględnia notowania, więc wynik jest optymistyczny.

Załóżmy teraz, dla przykładu, że książka ma dziesięć rozdziałów, a notowanie zajmuje pięć minut na rozdział. To pięćdziesiąt minut dodatkowo na każdą książkę. Przy trzydziestu minutach czytania dziennie, czyli 10\,950 minutach w roku, wychodzi:
\begin{przyklad}[Ile książek mieści się w roku]
Bez notowania: 10\,950 / 320 = 34,2 książki (320 minut to czas czytania jednej typowej książki).

Z notowaniem: 10\,950 / (320 + 50) = 10\,950 / 370 = 29,6 książki.

Różnica to około 4,6 książki, czyli około 13,5 procent. Dziesięć rozdziałów i pięć minut na rozdział to moje założenia, a nie dane o czytelnikach. Podstaw własne.
\end{przyklad}

Ten rachunek nie mówi, że notować nie należy. Mówi, że notatka jest inwestycją, którą trzeba uzasadnić przeznaczeniem książki. Notowanie w każdej książce, bez względu na to, czy wrócisz do niej jeszcze raz, obniża liczbę przeczytanych tytułów i nie daje za to niczego w zamian. Dla lektur, które mają działać w pracy, notatka potrafi się zwrócić wielokrotnie. Dla powieści zwykle nie.

\begin{zasada}[Zasada wyboru]
Notuj tylko to, co zmieni twoje działanie albo będzie potrzebne do powrotu. Reszta zostaje w pamięci albo odchodzi, i to też jest wynik.
\end{zasada}

\section{Jedno zdanie i trzy pomysły do zastosowania}

Najtańsza forma notatki to jedno zdanie napisane zaraz po skończeniu książki, zanim zamkniesz ją na dobre. Nie streszczaj rozdziałów. Napisz, co autor twierdzi, dlaczego tak twierdzi i co z tego wynika dla ciebie. Wzorzec jest prosty: ,,Autor twierdzi, że \dots, bo \dots, więc \dots''. Jeśli nie potrafisz wypełnić trzeciej części, nie wiesz jeszcze, czy książka jest dla ciebie. To także jest informacja.

Jedno zdanie samo nie wystarcza, bo wiedza bez działania szybko staje się opinią. Dlatego dołącz do niego trzy pomysły do zastosowania. Każdy powinien być konkretny: co zrobisz, kiedy i gdzie. ,,Będę bardziej uważny na sen'' nie jest pomysłem, bo nie da się tego sprawdzić. ,,W czwartek o 21:00 przesunę budzik na 6:30 i zapiszę, o której godzinie zasnąłem'' jest pomysłem, który można ocenić po tygodniu.

Pułapka jest tu jedna: trzy pomysły łatwo wypisać z książki, a trudno je przenieść do dnia. Przy każdym zapytaj, czy znasz pierwszy krok, termin i sposób sprawdzenia. Jeśli nie, pomysł zostaje na liście i nie zmienia niczego.

\section{Powtórki w odstępach i przypominanie przez rozmowę}

Jeśli chcesz, żeby treść została, musisz do niej wracać, a sposób powrotu ma znaczenie. Czytanie tej samej strony drugi raz daje poczucie znajomości, które bywa mylące: rozpoznajesz zdania, więc wydaje ci się, że je rozumiesz. Dużo skuteczniejsze jest \emph{przypomnienie}: zamknij książkę i spróbuj odtworzyć jej treść z pamięci, a dopiero potem sprawdź, co zostało pominięte.

Druga zasada to odstępy. Powtórka zaraz po lekturze jest tania, ale mało zapamiętasz, bo materiał jest jeszcze świeży i wydaje się oczywisty. Powtórka po kilku dniach kosztuje więcej wysiłku, ale to właśnie wysiłek wzmacnia ślad. Możliwy harmonogram dla jednej książki wygląda tak: dzień zero czytanie i jedno zdanie, dzień pierwszy trzy pomysły z pamięci, dzień siódmy krótkie przypomnienie kluczowego argumentu, dzień trzydziesty sprawdzenie, czy któryś pomysł został zastosowany. Taki harmonogram nie wymaga dużo czasu, około pięciu do dziesięciu minut w każdym z tych momentów.

Efekty odstępów i przypominania są dobrze opisane w badaniach nad uczeniem się, ale wielkość efektu w warunkach czytania dorosłych książek nie jest dla mnie pewna. \verify{korzyść z powtórek w odstępach jest realna w eksperymentach z materiałem do zapamiętania, a jej wielkość w czytaniu nieliterackim wymaga osobnego sprawdzenia w literaturze} Dlatego traktuj ten harmonogram jako eksperyment na sobie, a nie jako pewnik.

Przypominanie przez rozmowę ma jeszcze jedną zaletę. Jeśli musisz wyjaśnić komuś, o czym była książka, w zrozumiały sposób, nie możesz się ukryć za zdaniami autora. Wyjaśnienie wymaga przeformułowania, a przeformułowanie ujawnia, co naprawdę zrozumiałeś. Możesz to zrobić z przyjacielem, w klubie czytelniczym, a nawet z nagraniem głosowym, którego nikt nie usłyszy. Zasada jest taka: to, czego nie umiesz wyjaśnić, tego jeszcze nie rozumiesz.

\section{Zakreślanie, marginalia i cytaty}

Zakreślanie jest najpopularniejszą formą notowania i najmniej przydatną, jeśli jest jedyną. Kolorowy marker daje satysfakcję, bo zostawia ślad na stronie, ale nie wymaga od ciebie żadnego myślenia. Po miesiącu oglądasz strony pełne żółtych linii i nie wiesz, dlaczego każda z nich została zaznaczona. Zakreślenie jest zapisem uwagi w chwili czytania, a nie wnioskiem.

Reguła, która działa: zakreślaj rzadko, a zakreślając, dopisz jedno zdanie na marginesie. Co najwyżej jedno zakreślenie na stronę. Marginalia powinny odpowiadać na konkretne pytanie: ,,Gdzie to stosuję?'', ,,Czy to się zgadza z tym, co czytałem wcześniej?'', ,,Co by się stało, gdyby tu było inaczej?''. Pytanie zmusza do przetworzenia i tworzy coś, do czego możesz wrócić.

Cytaty są najbardziej kuszące i najłatwiej je przeładować. Zapisuj je z numerem strony, jednym zdaniem, bez całych akapitów. Przy każdym cytacie dopisz własnymi słowami, dlaczego go zapisałeś. Bez tego zostaje tekst bez kontekstu, który za rok będzie wyglądał na ciekawy, a nie będzie mówił, czemu miał być ważny. Limit, który sprawdza się u wielu osób, to kilka cytatów na książkę. Jeśli zapisujesz ich dwadzieścia, to notujesz, a nie czytasz.

\section{Dziennik czytania i łączenie książek}

Dziennik czytania to jedno miejsce, w którym zbierasz każdą książkę. Każdy wpis zawiera datę, tytuł, jedno zdanie streszczenia, jeden pomysł, który zastosowałeś lub zamierzasz zastosować, i jedno zdanie o tym, komu książkę poleciłbyś i dlaczego. Szablon i sposób prowadzenia dziennika opisuje Dodatek C. Dziennik jest najtańszą formą notatek, jaką znam, i zarazem najczęściej porzucaną, bo nie daje natychmiastowej nagrody. Jego wartość pojawia się po kilkunastu wpisach, gdy możesz przejrzeć, co naprawdę przeczytałeś w roku.

Druga korzyść jest mniej oczywista: łączenie książek. Książki rzadko mówią tę samą rzecz, a często mówią o tym samym problemie innymi słowami. Gdy dwie lektury opisują ten sam nawyk z różnych stron, ich zestawienie daje więcej niż każda z osobna. Zapisuj takie powiązania w dzienniku jako krótkie zdanie z numerami wpisów: ,,to samo co w wpisie 12, ale z innej perspektywy''. Zapisuj także sprzeczności. Jeśli jedna książka twierdzi, że powinieneś czytać rano, a druga, że wieczorem, nie wybieraj jednej z nich. Zanotuj, w jakich warunkach każda może mieć rację.

Połączenia pojawiają się w czytaniu dopiero wtedy, gdy masz gdzie je umieścić. Bez dziennika każda książka zostaje w osobnej szufladzie. Z dziennikiem zaczynasz zauważać wzorce: które tematy wracają, które pomysły się powtarzają i które książki warto przeczytać razem.

\begin{falle}[Pułapka]
Najczęstszy sposób, w jaki notowanie psuje czytanie, to przesunięcie uwagi. Zaczynasz od przyjemności lektury, po chwili zauważasz, że lepiej byłoby zapisać coś, co właśnie przeczytałeś, i nagle większość sesji poświęcasz na pisanie. Książka czyta się wolniej, bo co kilka zdań przerywasz, żeby zanotować. Liczba notatek rośnie, a liczba skończonych książek stoi w miejscu.

Objaw jest prosty do sprawdzenia: porównaj, ile stron zapisałeś w miesiącu z liczbą przeczytanych stron. Jeśli notatek jest więcej niż zdań przeczytanych przez ciebie w tym samym czasie, notowanie przejęło czytanie. Korekta ma charakter organizacyjny, a nie moralny. Ustal limit czasu na notatki w sesji, na przykład jedną piątą czasu czytania, i notuj dopiero po zakończeniu rozdziału. Notatka pisana w trakcie zdania zwykle jest tylko powtórzeniem tego zdania.
\end{falle}

\section{Tydzień bez zbierania}

Ostatnia część jest krótka, bo tu zwykle kończy się teoria. Notowanie jest narzędziem, a narzędzie dobiera się do celu. Z tego, co napisano wyżej, wynika prosty plan: najpierw odpowiedz na trzy pytania o użyteczność, potem wybierz najwyżej jedną formę notowania dla danej książki, a na końcu sprawdź po miesiącu, czy któryś z pomysłów faktycznie zmienił twoje działanie.

Dobrą miarą jest stosunek. Ile razy w miesiącu zrobiłeś coś inaczej dzięki książce? Jeśli odpowiedź brzmi ,,nigdy'', problem nie leży w notatkach. Leży w tym, że książki nie mają miejsca w dniu, albo że pomysły nie mają konkretnego pierwszego kroku. Notatki nie naprawią tego, czego nie zaplanowałeś.

\begin{cwiczenia}
\exitem{Wybierz trzy książki, które przeczytałeś w ostatnim roku. Dla każdej odpowiedz na trzy pytania: czy będę z niej korzystać, czy będę do niej wracać, czy sama lektura ma wartość. Zdecyduj, czy notować.}{Odpowiedź zależy od czytelnika, ale wzorcowy wynik to jasna decyzja dla każdej książki. Książka, na którą odpowiadasz ,,nie'' na wszystkie trzy pytania, nie wymaga notatek. Książka z jednym ,,tak'' dostaje jedną, krótką notatkę: jedno zdanie i jeden pomysł.}
\exitem{Oblicz, ile książek mieści się w roku przy trzydziestu minutach czytania dziennie, jeśli każda książka ma osiemdziesiąt tysięcy słów, czytasz z szybkością 250 słów na minutę i notujesz przez dwadzieścia minut na książkę.}{80\,000 / 250 = 320 minut na książkę bez notatek. Z notatkami: 320 + 20 = 340 minut. Rocznie: 30 \cdot 365 = 10\,950 minut. Bez notatek: 10\,950 / 320 = 34,2 książki. Z notatkami: 10\,950 / 340 = 32,2 książki. Różnica to około dwóch książek w roku.}
\exitem{Napisz jednozdaniowe streszczenie książki w formie ,,Autor twierdzi, że \dots, bo \dots, więc \dots''.}{Wzorzec wymaga trzech części. Przykład: ,,Autor twierdzi, że nawyk trzyma się lepiej, gdy jest przywiązany do istniejącej czynności, bo przypomnienie pojawia się samo, więc warto wybrać czynność, która zdarza się codziennie.'' Jeśli nie umiesz wypełnić trzeciej części, zapisz to jako wniosek do sprawdzenia.}
\exitem{Zamień poniższy pomysł na konkretne działanie: ,,Będę więcej czytać.''}{Konkretne działanie zawiera czas, miejsce i sprawdzenie. Przykład: ,,Od poniedziałku po obiedzie, w kuchni, przez dwadzieścia minut; po tygodniu policzę strony i zapiszę je w dzienniku''. Pomysł, którego nie da się ocenić po tygodniu, nie jest jeszcze pomysłem.}
\exitem{Zaplanuj powtórki dla książki przeczytanej dziś. Podaj dni, na które je zaplanujesz, i co dokładnie będziesz robić w każdej z nich.}{Przykładowy harmonogram: dzień zero, jedno zdanie i trzy pomysły; dzień pierwszy, przypomnienie trzech pomysłów z pamięci bez zaglądania do książki; dzień siódmy, przypomnienie głównego argumentu i sprawdzenie, czy pomysł został zrobiony; dzień trzydziesty, ocena efektu. Każda z tych czynności zajmuje około pięciu do dziesięciu minut.}
\exitem{Przeczytaj jeden rozdział i zrób tylko jedną notatkę: zamiast zakreślać, napisz pytanie na marginesie. Co to pytanie zmienia w twoim czytaniu następnego rozdziału?}{Odpowiedź zależy od tekstu. Wzorcowy efekt jest taki, że pytanie zmusza do zatrzymania się i przypomnienia sobie treści, a więc do wysiłku, który wzmacnia zapamiętanie. Jeśli po przeczytaniu następnego rozdziału nie pamiętasz pytania, pytanie było za ogólne.}
\exitem{Co zrobić, gdy dwie książki mówią sprzeczne rzeczy na ten sam temat?}{Nie wybieraj jednej z nich. Zanotuj w dzienniku, w jakich warunkach każda z nich ma rację, na przykład przy jakim typie pracy, poziomie doświadczenia albo skali problemu. Taka notatka zachowuje spór i daje ci kryterium, które pozwoli rozstrzygnąć go w konkretnej sytuacji.}
\end{cwiczenia}

\zadanie{Wybierz książkę, którą czytasz teraz. Przez tydzień notuj wyłącznie jedno zdanie po każdej sesji i nie dłużej niż jedną piątą czasu czytania. W dniu siódmym zamknij książkę i w ciągu pięciu minut zapisz z pamięci jej główny argument oraz trzy pomysły. Miara sukcesu: co najmniej dwa pomysły są poprawnie odtworzone, a przy każdym z nich wpisałeś pierwszy krok z terminem. Zapisz wynik w dzienniku czytania.}
